\section{Introduction}\label{sec:intro}

Camera-equipped smart glasses are becoming everyday devices, and an assistant on them is most useful when the wearer can refer to things in view. Pointing is the most direct way to do so: ``what is \emph{that}?'', ``remind me to buy \emph{this}''. Systems have combined speech with pointing since Bolt's \emph{Put-that-there}~\cite{bolt1980put}, and recent wearable assistants use pointing and gaze to resolve such pronouns~\cite{lee2024gazepointar}. Before any of this works on glasses, the device must answer a spatial question from its own camera: \emph{which object is the wearer pointing at?}

Answering this on glasses is hard for three reasons.
\begin{itemize}
\item \emph{The egocentric view.} The camera sees the wearer's hand from behind and often only in part. The target may be touched or across the room, and a pointing direction drawn in the image is only a rough guide to it.
\item \emph{Other people's hands.} A companion's hand or a stranger's pointing gesture can be in view, and must not be mistaken for the wearer's.
\item \emph{The budget.} Glasses-class processors share a battery of a few hundred milliampere-hours with the camera, audio and radio. Sending frames to a cloud model adds network latency and moves imagery of bystanders off the device.
\end{itemize}

\paragraph{Prior work.} Pointing-target grounding in egocentric images is an active research area, and EgoPoint-Ground~\cite{egopoint} provides a large benchmark for it. The strongest methods are large models:
\begin{itemize}
\item multimodal LLMs that reason step by step from the hand to a ray and then to the target~\cite{egopoint,li2026pointvgr};
\item a 226M-parameter transformer that models the pointing ray explicitly~\cite{li2026vistaref}.
\end{itemize}
Some benchmark tasks even give the model the ground-truth hand box~\cite{egopoint}. These systems reach high accuracy on a GPU, but none of them is designed for glasses-class hardware, and to our knowledge none addresses other people's hands in view. Egocentric fingertip systems, including one that runs on a phone, stop at the fingertip~\cite{huang2016pointing,jain2019gestarlite} and never localise the target.

\paragraph{SpatialNet.} We ask how far a model that fits on the glasses can go. SpatialNet needs only one 2D RGB frame:
\begin{itemize}
\item no text prompt, no hand box, no depth;
\item no external hand detector or pose estimator.
\end{itemize}
It predicts the wearer's hand, the fingertip, the pointing direction and the box of the pointed-at object, with presence gates for each. All outputs come from one forward pass of a compact convolutional network with no self-attention and no token-by-token decoding: 0.72M parameters, 0.23 GFLOPs, 1.6\,MB in FP16. This makes it 300 times smaller than VistaRef and about 10{,}000 times smaller than the 7B multimodal LLMs.

One design choice lets such a small network localise precisely: \emph{instance-anchored heads}. The heads form a chain in which each one reads a shared feature map at the location the previous one predicted, so the fingertip, direction and target all belong to the same hand.

We also give the target head the geometry of the predicted ray, as six ray-coordinate maps entering through a zero-initialised branch. We then test whether a model this small makes use of it, rather than assuming so. Training additionally pastes unlabelled hands of other people into the frames and penalises responses to them, so the model learns to prefer the wearer's hand without labelled bystander hands.

\paragraph{Evaluation.} We train and test on EgoPoint and report:
\begin{itemize}
\item accuracy by fingertip--target distance, which matters because about 70\% of targets are touched by the fingertip;
\item comparisons with published results of state-of-the-art models and zero-shot multimodal LLMs on the same test split, and with simple fingertip-based references;
\item comparisons with compact detectors (YOLO26n and NanoDet-Plus-m) that we trained on the same split, at the same compute, and scored with the same code;
\item ablations of each component;
\item behaviour when other people's hands are in view (pasted and real);
\item latency, memory and battery drain on an engineering unit of a pair of smart glasses.
\end{itemize}
\paragraph{Results in brief.}
\begin{itemize}
\item SpatialNet finds the wearer's hand reliably (P@0.5 of 92.6\%, the share of test frames with IoU${\geq}0.5$) and the fingertip precisely (3.9\,px mean error at $160\times160$).
\item It localises the pointed-at object with a P@0.5 of 58.9\% (mIoU 53.3\%). That is comparable to zero-shot 7--8B multimodal LLMs reported on the same test split (44.5--49.8\%), close to a 7B model fine-tuned on EgoPoint (63.6\%), and far below VistaRef (93.9\%) and PointVG-R (83.6\%).
\item At the same compute, compact detectors with 1.6--3.5 times the parameters draw tighter boxes (P@0.5 59.9--67.6\%), but SpatialNet places the target centre on the right object at least as often (84.8\% against 81.6--82.9\%), and only SpatialNet keeps its accuracy when other hands are pasted into view.
\item It drops to 25.4\% for far targets. Trained at higher input resolution, with no new parameters, the same model reaches 62.4\% at 224\,px (far: 35.9\%) and 65.1\% at 320\,px (Section~\ref{sec:resolution}).
\item It runs in 60--80\,ms per frame on the glasses. Running it continuously costs less battery (9\% per minute on the test unit) than streaming frames to a phone at 0.5\,fps (16\%).
\end{itemize}
We present this as an operating point, with its limits measured, not as a replacement for large models.

\paragraph{Contributions.}
\begin{enumerate}
\item \emph{An image-only, single-pass pointed-target localiser that runs on smart glasses.}
  \begin{itemize}
  \item From one RGB frame, a 0.72M-parameter CNN without self-attention, external models or prompts outputs the hand, fingertip, direction and target box.
  \item On a smart-glasses engineering unit it takes 60--80\,ms per frame, and running it continuously costs less battery than streaming frames off the device (Sections~\ref{sec:method}, \ref{sec:device} and~\ref{sec:ondevice}).
  \end{itemize}
\item \emph{Instance-anchored heads that keep a sub-million-parameter model precise.} Against a plain model of the same size, the anchoring improves:
  \begin{itemize}
  \item the hand box, by 1.8 points of mIoU;
  \item the direction, by 3.4$^\circ$;
  \item P@0.7, by 2.8 points (1.7 on the test split without duplicates, where the interval reaches zero);
  \item small targets, by 11.3 points.
  \end{itemize}
  The first two gains are well beyond the variation between training seeds; the last two are supporting evidence (Sections~\ref{sec:heads} and~\ref{sec:ablations}).
\item \emph{An analysis of what limits pointing localisation at this scale.}
  \begin{itemize}
  \item Input resolution is the largest lever, at no parameter cost.
  \item Distractor pasting makes the model robust to other hands at no cost on clean images.
  \item A prior learned from explicit ray geometry collapses into a fingertip-proximity prior when most targets are touched. This explains why far targets stay hard, and indicates how future models should constrain the ray (Sections~\ref{sec:resolution}--\ref{sec:own}).
  \end{itemize}
\item \emph{An evaluation that places accuracy next to cost.} It compares with published models that are 300 to 10{,}000 times larger and with compact detectors re-trained at the same compute, stratifies by fingertip--target distance, and decomposes the errors (Section~\ref{sec:results}).
\end{enumerate}
Distractor pasting is treated as a supporting result rather than a headline contribution. Its effect is clear (Section~\ref{sec:ablations}), but its test distractors are synthetic.
